\documentclass[11pt]{article}
\usepackage{amsmath,amssymb}
\usepackage{physics}
\usepackage[margin=1in]{geometry}
\usepackage{enumitem}

\title{ATPL 2025 --- Exercises: Measurement}
\author{Introduction to Quantum Computing}
\date{}

\begin{document}
\maketitle

Background reading: Nielsen \& Chuang, Section 2.2.3 (projective
measurements).

\begin{enumerate}[label=\textbf{Exercise \arabic*.}, leftmargin=2.2cm]

\item \textbf{Measurement probabilities in the computational basis.}
Let $\ket\psi = \sqrt{1/5}\,\ket0 + \sqrt{4/5}\,\ket1$.
\begin{enumerate}[label=(\alph*)]
\item What are the probabilities of obtaining $\ket0$ and $\ket1$ when
measuring $\ket\psi$ in the $\{\ket0,\ket1\}$ basis?
\item What is the post-measurement state in each case?
\end{enumerate}

\item \textbf{Measuring in a rotated basis.}
Using the same $\ket\psi$ as above, express it in the $\{\ket+,\ket-\}$
basis and compute the probabilities of obtaining $\ket+$ and $\ket-$.
Compare with Exercise 1: why can the same physical state give very
different-looking measurement statistics depending on the chosen basis?

\item \textbf{A state that is deterministic in one basis, random in
another.}
Take $\ket\psi = \ket+$.
\begin{enumerate}[label=(\alph*)]
\item What is the outcome probability distribution when measuring in the
$\{\ket+,\ket-\}$ basis?
\item What is the outcome probability distribution when measuring in the
$\{\ket0,\ket1\}$ basis? Show your work using
$\ket+=\tfrac1{\sqrt2}(\ket0+\ket1)$.
\end{enumerate}

\item \textbf{Partial measurement of an entangled state.}
Consider $\ket\psi = \tfrac1{\sqrt2}(\ket{00}+\ket{11})$.
\begin{enumerate}[label=(\alph*)]
\item If the first qubit is measured in the $\{\ket0,\ket1\}$ basis,
list the possible outcomes, their probabilities, and the resulting
2-qubit post-measurement state in each case.
\item Suppose instead you measure only the \emph{second} qubit. Argue
(by symmetry of $\ket\psi$, or by direct computation) that you get the
same outcome probabilities and correlations as in part (a).
\item Now consider $\ket\varphi = \tfrac12(\ket{00}+\ket{01}+\ket{10}-\ket{11})$.
Measure the first qubit in the computational basis: what are the
probabilities and post-measurement states? Is the resulting second-qubit
state, conditioned on the outcome, itself a definite classical bit
value, or a superposition?
\end{enumerate}

\item \textbf{Unequal superpositions and biased coins.}
A biased "quantum coin" is prepared as
$\ket\psi = \sqrt p\,\ket0 + \sqrt{1-p}\,\ket1$ for $p\in[0,1]$ (ignoring
relative phase for simplicity).
\begin{enumerate}[label=(\alph*)]
\item For what value of $p$ does measuring $\ket\psi$ in the
computational basis reproduce a fair classical coin flip?
\item Design (in words, or as a short gate sequence using an
$R_Y(\theta)$ rotation of your choice) a way to prepare $\ket\psi$ for
general $p$ starting from $\ket0$.
\end{enumerate}

\item \textbf{Measurement destroys information: no-cloning intuition.}
Explain in your own words, using the results from Exercises 1--2, why
you cannot recover the full description $(a_0,a_1)$ of an unknown qubit
state $\ket\psi=a_0\ket0+a_1\ket1$ from a single measurement, or even
from many measurements of a \emph{single} copy of $\ket\psi$. What
\emph{can} you learn about $(a_0,a_1)$ if you are given many independent
identically prepared copies of $\ket\psi$ and are free to choose your
measurement basis?

\item \textbf{Projective measurement as a projector (optional, more
formal).}
A projective measurement of qubit $k$ with outcome $v\in\{0,1\}$ can be
written as applying the projector $P_v^{(k)} = I^{\otimes(k-1)}\otimes
\ket v\!\bra v \otimes I^{\otimes(n-k)}$ to an $n$-qubit state and then
renormalizing.
\begin{enumerate}[label=(\alph*)]
\item Verify that $P_0^{(1)}$ and $P_1^{(1)}$ applied to
$\ket{\Phi^+}=\tfrac1{\sqrt2}(\ket{00}+\ket{11})$ reproduce the two
post-measurement states you found in Exercise 4(a), including correct
normalization.
\item Show that $P_0^{(k)}+P_1^{(k)} = I$, and explain why this is
exactly the statement that "the probability of measuring something is
always 1" from the lecture.
\end{enumerate}

\end{enumerate}

\end{document}
